Nullniveau der Lageenergie: die Höhe des Schwerpunkts beim Anlauf. Zustand 1: beim Anlauf (nur kinetische Energie). Zustand 2: im höchsten Punkt (nur Lageenergie; nimm an, sie sei dort in Ruhe).

\begin{abcliste}
\abc
Sei $v$ ihre Geschwindigkeit, $h_0$ die Höhe ihres Schwerpunkts beim Anlauf und $\Delta h$, um wie viel er steigt. Die kinetische Energie im Zustand 1 wird zur Lageenergie im Zustand 2:
\begin{align}
 \frac{1}{2}\,m\,v^2 &= m\,g\,\Delta h
\end{align}
Der Schwerpunkt steigt von seiner Höhe beim Anlauf aus:
\begin{align}
 \ssc{h}{max} &= \hmaxF \\
   &= \hA+\frac{\left(\vA\right)^2}{2\cdot\g} \\
   &= \hmax \approx \result{\hmaxP{0}{3}}
\end{align}

\abc
Sei $H$ die Höhe, die der Schwerpunkt erreichen soll. Nach der Geschwindigkeit aufgelöst:
\begin{align}
 v &= \vnF \\
   &= \sqrt{2\cdot\g\cdot(\hZ-\hA)} \\
   &= \vn \approx \result{\vnP{0}{3}}
\end{align}
Spitzenathletinnen sprinten tatsächlich mit etwa \SI{10}{\mps}. Sie stossen sich zudem mit den Armen ab, und ihr Schwerpunkt passiert knapp unter der Latte.
\end{abcliste}
